\documentclass[11pt]{article}
\usepackage{mathunal-base}
\mathunalfuente{serif}

\renewcommand{\examtitulo}{Cálculo en Varias Variables --- Segundo Parcial (Tema A)}
\renewcommand{\examfecha}{Semestre 01-2023, 17 de abril de 2023}

\begin{document}

\examheader

\vspace{4pt}
\noindent Nombre: \dotfill\ Cédula: \dotfill

\vspace{4pt}
\noindent Grupo: \dotfill\ Firma: \dotfill

\vspace{4pt}
\noindent Puntaje. Solo para uso oficial:\quad 1--3 (/40) \dotfill;\quad 4 (/20) \dotfill;\quad 5 (/20) \dotfill;\quad 6 (/20) \dotfill.\quad Total: \dotfill\quad Nota: \dotfill

\vspace{6pt}
\noindent\textbf{Instrucciones.} La duración del examen es de 1 hora y 50 minutos. El examen consta de seis puntos en dos hojas impresas por ambos lados; verifique que su examen esté completo. En las preguntas con procedimiento justifique sus respuestas en los espacios asignados. No está permitido sacar hojas en blanco ni ningún tipo de apuntes durante el examen; verifique que su celular esté apagado. No se permite el uso de calculadora.

\vspace{8pt}
\noindent\textbf{En las preguntas 1 a 3 no se tendrá en cuenta el procedimiento.}

\begin{enumerate}
\item{} (10\,\%) Sea $f: \mathbb{R}^2 \to \mathbb{R}$ una función continua. Exprese la integral
\[
\int_0^1 \int_{y^2}^{\sqrt y} f(x,y)\,dx\,dy
\]
como una integral definida en la forma $\displaystyle\iint_D f(x,y)\,dy\,dx$.

\vspace{2pt}
\noindent Respuesta: \underline{\phantom{XXXXXXXXXXXXXXXX}}\quad (Escriba claramente cada límite de integración.)

\item{} (20\,\%) Considere el sólido en $\mathbb{R}^3$ al interior del cilindro $x^2 + y^2 = 9$, entre los planos $z = 1$ y $z = 5$. Y suponga que la densidad en cada punto $(x,y,z)$ es igual a la distancia del punto al eje $y$.
\begin{itemize}
\item[(a)] (10\,\%) Escriba una integral definida en coordenadas cartesianas (rectangulares), mediante la cual se pueda calcular la masa del sólido (no la evalúe).

\vspace{2pt}
\noindent Respuesta: \underline{\phantom{XXXXXXXXXXXXXXXX}}
\item[(b)] (10\,\%) Escriba la misma integral del ítem (a) en coordenadas cilíndricas (no la evalúe).

\vspace{2pt}
\noindent Respuesta: \underline{\phantom{XXXXXXXXXXXXXXXX}}
\end{itemize}

\item{} (10\,\%) Plantee una integral en coordenadas esféricas que permita calcular la masa del sólido, en el primer octante, limitado por la esfera $x^2 + y^2 + z^2 = 16$, y cuya densidad en cada punto $(x,y,z)$ es inversamente proporcional al cuadrado de la distancia del punto al origen (no la evalúe).

\vspace{2pt}
\noindent Respuesta: \underline{\phantom{XXXXXXXXXXXXXXXX}}
\end{enumerate}

\vspace{4pt}
\noindent\textbf{En los ejercicios 4, 5 y 6 debe realizar el procedimiento justificando cada paso.}

\begin{enumerate}[start=4]
\item{} (20\,\%) Sea $f(x,y) = (x^3 - x)(y - 1)$. Halle los puntos críticos de $f$ y clasifíquelos usando el criterio de las segundas derivadas.

\item{} (20\,\%) Determine el máximo y el mínimo de la función $f(x,y) = 9x^2 + y^2 + xy$ para los puntos $(x,y)$ que satisfacen $9x^2 + y^2 = 4$.

\item{} (20\,\%) Encuentre el momento respecto al eje $x$, $M_x$, de la región del primer cuadrante limitada por $y = x$, $y = \sqrt3\,x$ y $x^2 + y^2 = 9$, y cuya densidad en cada punto es $\rho(x,y) = \cos\!\left((x^2+y^2)^{3/2}\right)$.
\end{enumerate}

\examfooter

\end{document}
